% =======================================================================================
% MODELO DE TCC - INSTITUTO FEDERAL DE SERGIPE (IFS)
% =======================================================================================
% Autores do Modelo (discentes de Ciência da Computação - IFS, Campus Itabaiana):
%   Mikael dos Santos Gois
%   Talyson Rodrigues Silva Nascimento
%   Kauan Gutierrez Santos Estácio
% Modelo Base: abnTeX2, adaptado para as normativas internas de TCC do IFS (Art. 31 e 34)
% Créditos dos modelos base (abnTeX2, DECOM/UFOP e Diogo D. Moreira) e das imagens: ver NOTICE
% Licença: CC BY-SA 4.0 (ver LICENSE)
% Este arquivo faz parte de um modelo opensource para estudantes do IFS.
% Não remova os créditos dos autores originais e do modelo base.
% Este modelo é fornecido "no estado em que se encontra", sem garantias de qualquer tipo.
% =======================================================================================

% [O que aparece no PDF Bookmark] {O que aparece impresso no documento}
\chapter[Metodologia]{METODOLOGIA} \label{sec:metodologia}

% ------------------------------------------------------------------------------
% EXEMPLO DE ALGORITMO / CÓDIGO-FONTE
% ------------------------------------------------------------------------------
% DICA PARA O ESTUDANTE:
% - O comando \ref{} referencia o número do algoritmo gerado automaticamente.
% - Referencie o algoritmo no texto ANTES dele aparecer, como na frase abaixo.
% - O código-fonte é renderizado em PRETO E BRANCO (NBR 14724:2024, 5.1), com
%   destaque de palavras-chave por negrito/itálico, e não por cor.
% - Mantenha o código alinhado à esquerda (sem espaços antes da 1ª coluna),
%   preservando apenas a indentação interna da linguagem.
% - SUBSTITUA este exemplo pelo seu próprio algoritmo antes de entregar o
%   trabalho. Para trocar a linguagem, mude "language=Java".
\section{Exemplo de Inserção de Algoritmo}

Para demonstrar a inserção de um algoritmo no texto, apresenta-se a seguir um exemplo de código em Java no Algoritmo \ref{alg:exemplo_java}. Note que o algoritmo é citado antes de sua apresentação e que o bloco é gerado automaticamente pelo ambiente de algoritmos.
\begin{algorithm}[htb]
    \caption{Exemplo de Algoritmo em Java}
    \label{alg:exemplo_java}
    \begin{lstlisting}[language=Java]
public static int somar(int a, int b) {
    // Retorna a soma de dois inteiros
    return a + b;
}
    \end{lstlisting}
\end{algorithm}

% DICA: Lembre-se de citar figuras, tabelas e algoritmos no texto com \ref{} e
% de inserir a fonte das ilustrações com \fonte{}.